%!TEX root = ../main.tex

%--------------------------------------------------------%
% DOCUMENT CLASS
%--------------------------------------------------------%

  \documentclass[letterpaper,12pt]{article}

%--------------------------------------------------------%
% TITLE SECTION
%--------------------------------------------------------%

  %Abstract
  \usepackage{abstract}
  \renewcommand{\abstractnamefont}{\normalfont\bfseries}
  \renewcommand{\abstracttextfont}{\normalfont\small\itshape}

  %Title
  \usepackage{titlesec}

  %Authors
  \usepackage{authblk}

  %Date
  \usepackage{datetime2}
  \DTMsetstyle{mdyy} % US date format

%--------------------------------------------------------%
% HEADERS & FOOTERS
%--------------------------------------------------------%

  \pagestyle{plain}

%--------------------------------------------------------%
% MACROS
%--------------------------------------------------------%

  \providecommand{\keywords}[1]{\textbf{\textit{Keywords---}} #1}

%--------------------------------------------------------%
% MATH SUPPORT
%--------------------------------------------------------%

  \usepackage{amssymb}
  \usepackage{amsthm}
  \usepackage{newtxmath}
  \usepackage{mathtools}
  \usepackage{blkarray, bigstrut}

%--------------------------------------------------------%
% FONTS
%--------------------------------------------------------%

  \usepackage{microtype}
  \usepackage{newtxtext}

%--------------------------------------------------------%
% LINES
%--------------------------------------------------------%

  \usepackage{setspace}
  \usepackage{lineno}
  \usepackage[table,x11names]{xcolor}
  \usepackage{enumitem}
  \setlist{nosep}
  \setlist[itemize]{leftmargin=*}

%--------------------------------------------------------%
% MARGINS
%--------------------------------------------------------%

  \usepackage[top=1in, bottom=1in, left=1in, right=1in]{geometry}

%--------------------------------------------------------%
% COMMENTS
%--------------------------------------------------------%

  \usepackage[colorinlistoftodos]{todonotes}
  \usepackage{soul}
  \usepackage{comment}
  \usepackage{xcolor}
  \usepackage{xcolor}
  \usepackage{framed}

  \newif\ifshowalts
  \showaltstrue

  \newenvironment{alttext}[1][]{
    \ifshowalts
      \begin{framed}
      \noindent\textbf{ALT TEXT}\ifx\\#1\\\else\ \textbf{(#1)}\fi\par
      \vspace{0.5ex}
    \else
      \expandafter\comment
    \fi
  }{
    \ifshowalts
      \end{framed}
    \else
      \endcomment
    \fi
  }
  \newcommand{\altsentence}[2]{%
  \ifshowalts
    \begin{alttext}[#1]
    #2
    \end{alttext}
  \fi}

%--------------------------------------------------------%
% ACRONYMS
%--------------------------------------------------------%

  \usepackage[nohyperlinks,nolist]{acronym}

%--------------------------------------------------------%
% GRAPHICS
%--------------------------------------------------------%

  \usepackage{graphicx}
  \usepackage{caption}
  \usepackage{subcaption}
  \graphicspath{{figures/}}
  \usepackage{float}
  \usepackage{printlen}
  \usepackage[section]{placeins}

%--------------------------------------------------------%
% TABLES
%--------------------------------------------------------%

  \usepackage{longtable}
  \newcommand{\sym}[1]{\rlap{#1}}
  \usepackage{tabularx}
  \newcolumntype{L}[1]{>{\raggedright\arraybackslash}p{#1}}
  \newcolumntype{C}[1]{>{\centering\arraybackslash}p{#1}}
  \newcolumntype{R}[1]{>{\raggedleft\arraybackslash}p{#1}}
  \usepackage{makecell}
  \usepackage{relsize}
  \usepackage{multirow}
  \usepackage{pdflscape}

%--------------------------------------------------------%
% REFERENCES (load hyperref last!)
%--------------------------------------------------------%

  \usepackage{csquotes}
  \usepackage[
    style=nature,
    doi=true,
    url=false,
    backend=biber,
    sorting=none,
    maxbibnames=3,
    minbibnames=3,
  ]{biblatex}
  \DeclareFieldFormat{doi}{\url{https://doi.org/#1}}
  \bibliography{references.bib}
  % xr-hyper must precede hyperref; resolves \ref across main.tex <-> supplement.tex
  \usepackage{xr-hyper}
  \usepackage{hyperref}
  % let long URLs wrap at slashes, hyphens and dots instead of running into the margin
  \def\UrlBreaks{\do\/\do-\do.}
